\chapter{Curvature of Curves and Polyhedra}\label{geometry:ch:curvature}

\begin{center}\itshape
If a curved surface is developed upon any other surface whatever, the measure of curvature at each point remains unchanged.
\upshape(paraphrase of Gauss's \emph{Theorema Egregium}, 1827)
\end{center}

In 1827 Carl Friedrich Gauss, who had spent years surveying the Kingdom of Hanover with theodolites and chains, published a short memoir on curved
surfaces. Its central result, which he himself called \emph{egregium}, ``remarkable'', says that the curvature of a surface can be measured by a creature
living in it, who knows nothing of the space around: it depends only on lengths measured on the surface. One consequence is familiar to everyone who has
eaten a slice of pizza (bend it sideways, and it cannot droop) or looked at a world map (no flat map of the Earth can get all distances right).

Chapter~4 measured curvature from inside a surface, by lengths alone, and proved Gauss--Bonnet for a smooth surface. This appendix is the elementary account, and it is not a second proof of that theorem. It treats the curvature of a plane curve, the four-vertex theorem, and Descartes' theorem that the defects at the corners of a polyhedron add to $4\pi$. A reader who has not read Chapter~4 still gets those theorems. The same total, fixed by the topology, has already appeared in the angle sums of spherical triangles and of hyperbolic ones (Appendix~I).

\begin{think}
Take a sheet of paper and roll it into a cylinder: are lengths drawn on the sheet preserved? Now try to wrap the sheet around a ball without creasing it.
What goes wrong?
\end{think}

\section{The curvature of a plane curve}

Let $\gamma:[0,L]\to\mathbb R^2$ be a curve of class $C^2$ parametrised by arc length, so that $T(s)=\gamma'(s)$ is a unit vector. Write
$T(s)=(\cos\theta(s),\sin\theta(s))$, where the angle $\theta$ is chosen continuous; for instance $\theta(s)=\theta(0)+\int_0^s(x'y''-y'x'')\,du$ works, since
$x'y''-y'x''=\theta'$ for $(x',y')=(\cos\theta,\sin\theta)$. The \emph{curvature} is the speed at which the tangent turns:
\[
\kappa(s)=\theta'(s),\qquad\gamma''(s)=\kappa(s)N(s),
\]
where $N$ is $T$ rotated by $+90^\circ$. A circle of radius $r$ traversed counterclockwise has $\kappa=\frac1r$, a line has $\kappa=0$. For a curve with any regular
parametrisation $t\mapsto(x(t),y(t))$, the chain rule gives (Exercise~\ref{geometry:ex:graphcurv})
\[
\kappa=\frac{x'y''-y'x''}{(x'^2+y'^2)^{3/2}}.
\]
The \emph{osculating circle} at a point with $\kappa\ne0$ is the circle of radius $\frac1{|\kappa|}$ centred at $\gamma+\frac1\kappa N$: among all circles, it hugs the curve
best (Figure~\ref{geometry:fig:vertices}).

\begin{theorem}[curvature determines the curve]\label{geometry:thm:fundcurve}
For every continuous function $\kappa$ on $[0,L]$ there is a unit-speed curve with curvature $\kappa$, and two such curves differ by a rotation followed by a
translation.
\end{theorem}

\begin{proof}
Set $\theta(s)=\theta_0+\int_0^s\kappa$ and $\gamma(s)=\gamma_0+\int_0^s(\cos\theta,\sin\theta)$. Then $\gamma'=(\cos\theta,\sin\theta)$ is a unit vector and its angle
has derivative $\kappa$. Conversely, any unit-speed curve with curvature $\kappa$ has a continuous tangent angle with $\theta'=\kappa$, so $\theta=\theta_0+\int\kappa$, and
$\gamma=\gamma(0)+\int(\cos\theta,\sin\theta)$: changing $\theta_0$ rotates the curve, and changing $\gamma(0)$ translates it.
\end{proof}

For a closed curve, $\theta(L)-\theta(0)=\int_0^L\kappa\,ds$ is a multiple of $2\pi$, since the tangent comes back to itself: the \emph{total curvature} counts the
turns of the tangent. Heinz Hopf proved in 1935 that it is $\pm2\pi$ for a simple closed curve (a remark here; Exercise~\ref{geometry:ex:totalcurv} shows the figure eight,
with total curvature $0$).

\section{The four vertex theorem}

A \emph{vertex} of a curve is a point where the curvature has a local maximum or minimum. An ellipse has four, at the ends of its axes
(Figure~\ref{geometry:fig:vertices}). Could an egg-shaped curve have only two, the sharpest point and the flattest one?

\begin{figure}[ht]
\centering
\fig{geometry-ch12-vertices}
\caption{The ellipse $x=1.6\cos t$, $y=\sin t$ with its osculating circles at the four vertices (radius $0.625$ at the ends of the long axis, $2.56$ at the ends of
the short one), and its curvature as a function of $t$: two maxima and two minima.}
\label{geometry:fig:vertices}
\end{figure}

We call \emph{oval} a closed curve of class $C^3$, parametrised by arc length, that bounds a convex region and meets every line in at most two points.

\begin{theorem}[four vertex theorem]\label{geometry:thm:fourvertex}
The curvature of an oval that is not a circle has at least two local maxima and two local minima.
\end{theorem}

\begin{proof}
\emph{Three integrals vanish.} For every closed curve, with $\gamma=(x,y)$ and $(x',y')=(\cos\theta,\sin\theta)$,
\[
\int_0^L\kappa'\,ds=0,\qquad\int_0^L\kappa'x\,ds=-\int_0^L\kappa\cos\theta\,ds=-\int_0^L(\sin\theta)'\,ds=0,
\]
using integration by parts (all functions are $L$-periodic) and $\kappa\cos\theta=\theta'\cos\theta$; similarly $\int\kappa'y\,ds=0$. Hence
$\int_0^L\kappa'(s)\,\ell(\gamma(s))\,ds=0$ for every affine function $\ell(x,y)=ax+by+c$.

\emph{Two vertices are not enough.} The continuous function $\kappa$ has a maximum at some point $P$ and a minimum at some point $Q$, and $P\ne Q$ since $\kappa$ is
not constant (a curve of constant curvature $\frac1r$ is a circle, by Theorem~\ref{geometry:thm:fundcurve}). Suppose $\kappa$ has no other local extremum. Then $\kappa$ is
monotone on each of the two arcs from $P$ to $Q$: nonincreasing on one, nondecreasing on the other. Let $\ell$ be an affine function vanishing on the line $PQ$. This
line meets the oval only at $P$ and $Q$, so $\ell$ has a constant sign on each open arc; the signs are opposite, because the chord $PQ$ crosses the convex region
and the two arcs lie on the two sides of it. Choosing the sign of $\ell$, the product $\kappa'\,\ell(\gamma)$ is $\ge0$ everywhere. Its integral is $0$, so it vanishes
identically, so $\kappa'=0$ on both open arcs, and $\kappa$ is constant: a contradiction.

So $\kappa$ has a third local extremum. Going around the curve, strict local maxima and minima of $\kappa$ alternate (between two local maxima lies a minimum of $\kappa$ on
the arc joining them); if there are finitely many, their number is therefore even, hence at least four, two of each kind. If there are infinitely many, there is
nothing to prove.
\end{proof}

The theorem was proved for ovals by Syamadas Mukhopadhyaya in 1909, and for all simple closed curves by Adolf Kneser in 1912. The argument above, with the
affine function $\ell$, is due to G.~Herglotz. It explains why an egg always has four vertices: two pointed ends are not enough.

\section{The curvature of a surface}

Let $S\subset\mathbb R^3$ be a smooth surface and $p\in S$. After a rotation and a translation, $p=0$ and the tangent plane at $p$ is $z=0$, so that near $p$ the
surface is the graph of a smooth function $z=f(x,y)$ with $f(0)=0$ and $\nabla f(0)=0$. Its Taylor expansion starts with the quadratic form
$\frac12(x,y)\,H\,(x,y)^{\mathsf T}$, where $H$ is the Hessian of $f$ at $0$.

\begin{definition}
The \emph{principal curvatures} $\kappa_1,\kappa_2$ of $S$ at $p$ are the eigenvalues of $H$, and the \emph{Gaussian curvature} is $K=\kappa_1\kappa_2=\det H$.
\end{definition}

Changing the orientation of the normal changes the signs of $\kappa_1,\kappa_2$ but not $K$, and rotating the $(x,y)$-plane conjugates $H$ by a rotation, which does not
change its determinant; so $K$ is well defined. By the spectral theorem, $\kappa_1$ and $\kappa_2$ are the curvatures of the sections of $S$ by the two orthogonal vertical
planes through the eigenvectors. Examples: a sphere of radius $r$, near its south pole, is $z=r-\sqrt{r^2-x^2-y^2}=\frac{x^2+y^2}{2r}+\dots$, so $\kappa_1=\kappa_2=\frac1r$
and $K=\frac1{r^2}$. A cylinder of radius $r$ has $\kappa_1=\frac1r$, $\kappa_2=0$, $K=0$. The saddle $z=xy$ has $H=\begin{pmatrix}0&1\\1&0\end{pmatrix}$, so $K=-1$ at
the origin. Positive $K$: the surface lies on one side of its tangent plane, like an egg; negative $K$: it crosses it, like a saddle (Appendix~E).

The \emph{Gauss map} sends each point of $S$ to its unit normal vector $N(p)\in S^2$. For the graph,
$N=\frac{(-f_x,-f_y,1)}{\sqrt{1+f_x^2+f_y^2}}$.

\begin{proposition}\label{geometry:prop:gaussmap}
At $p$, the Jacobian determinant of the Gauss map, from the tangent plane of $S$ at $p$ to the tangent plane of $S^2$ at $N(p)$, equals $K(p)$. So for small regions
$U$ around $p$, $\operatorname{area}N(U)\approx|K(p)|\operatorname{area}(U)$.
\end{proposition}

\begin{proof}
At $0$, $\nabla f=0$, so $N(0)=(0,0,1)$ and both tangent planes are the plane $z=0$. Differentiating, the first two components of $N$ are $-f_x$ and $-f_y$ up to a
factor $1+O(|\nabla f|^2)$, whose derivative vanishes at $0$; so the differential of $N$ at $0$ is $-H$, with determinant $(-1)^2\det H=K$. The statement about
areas is the change of variables formula.
\end{proof}

\begin{remark}
Gauss's \emph{Theorema Egregium} says that $K$ can be computed from the lengths of curves on $S$ alone. So bending a surface without stretching it, as when a sheet of
paper is rolled into a cylinder, cannot change $K$. A plane has $K=0$ and a sphere of radius $r$ has $K=\frac1{r^2}$, so no piece of a sphere can be flattened without
distortion: every map of the Earth lies somewhere (Exercise~\ref{geometry:ex:nomap} proves this directly). And a pizza slice, curved downwards along one direction,
must stay straight in the other: $K=0$ forces one principal curvature to vanish. The proof of the theorem is a computation with the coefficients of the first
fundamental form; see do Carmo's book in the sources.
\end{remark}

\section{Gauss--Bonnet}

\begin{theorem}[Gauss--Bonnet for convex surfaces]\label{geometry:thm:gbconvex}
Let $S$ be the boundary of a compact convex body in $\mathbb R^3$. Assume that $S$ is a smooth surface with $K>0$ everywhere. Then
\[
\iint_SK\,dA=4\pi.
\]
\end{theorem}

\begin{proof}
We show that the Gauss map $N:S\to S^2$, with $N$ the outer normal, is a bijection. Given a unit vector $u$, the linear function $x\mapsto u\cdot x$ attains its maximum
on the compact set $S$ at some point $p$; there the plane $u\cdot x=u\cdot p$ supports the body, so it is the tangent plane at $p$ and $u=N(p)$. Conversely, if $u=N(p)$, the
tangent plane at $p$ supports the convex body, so $p$ maximises $u\cdot x$. Two points $p\ne q$ maximising $u\cdot x$ would give a segment $[p,q]$ in the support plane
and in the body, hence in $S$, along which $S$ contains a straight line, so that the curvature of $S$ in that direction vanishes and $K\le0$ there, which is excluded
(along the segment, the surface lies on one side of the tangent plane and contains a line of it, so $H$ is semidefinite with a zero direction). So $N$ is a bijection.

By Proposition~\ref{geometry:prop:gaussmap} its Jacobian is $K>0$, so $N$ is a local diffeomorphism, and being bijective it is a diffeomorphism. By the change of variables
formula, $\iint_SK\,dA=\operatorname{area}(N(S))=\operatorname{area}(S^2)=4\pi$.
\end{proof}

So the total curvature of every egg is $4\pi$, whatever its shape: flattening it somewhere forces it to be more curved elsewhere. Pierre Ossian Bonnet extended Gauss's
local formula in 1848, and the general theorem reads $\iint_SK\,dA=2\pi\chi(S)$ for every compact surface, where $\chi$ is the \emph{Euler characteristic}: $2$ for a sphere,
$0$ for a torus (Exercise~\ref{geometry:ex:torusK}), $2-2g$ for a surface with $g$ holes. Its local version, for a geodesic triangle $T$,
$\iint_TK\,dA=\alpha+\beta+\gamma-\pi$, contains Girard's theorem ($K=1$) and the hyperbolic area formula ($K=-1$) of Appendix~I.

\section{Polyhedra and Descartes' theorem}

The flat version of all this was found first. At a vertex $v$ of a convex polyhedron, the face angles add up to less than $2\pi$ (otherwise the vertex would be
flat); the \emph{angular defect} is $\delta_v=2\pi-(\text{sum of the face angles at }v)$. At a corner of a cube, three right angles meet, and $\delta_v=\frac\pi2$
(Figure~\ref{geometry:fig:defect}).

\begin{figure}[ht]
\centering
\fig{geometry-ch12-defect}
\caption{The three squares at a corner of a cube, unfolded into the plane, leave a gap of $\frac\pi2$: the angular defect. The cube has eight corners, and
$8\cdot\frac\pi2=4\pi$.}
\label{geometry:fig:defect}
\end{figure}

\begin{theorem}[Euler, 1758]\label{geometry:thm:euler}
A convex polyhedron with $V$ vertices, $E$ edges and $F$ faces satisfies $V-E+F=2$.
\end{theorem}

\begin{proof}
Project the boundary of the polyhedron radially onto a sphere centred at an interior point, and then stereographically from a point of the sphere inside one face
onto a plane. This gives a connected graph drawn in the plane without crossings, with $V$ vertices and $E$ edges, which cuts the plane into $F$ regions (the face
containing the pole becomes the unbounded region). We prove $V-E+F=2$ for every connected graph drawn in the plane, by induction on $E$. If the graph has no
cycle, it is a tree: $E=V-1$ and there is one region, so $V-E+F=V-(V-1)+1=2$. Otherwise remove an edge $e$ lying on a cycle. The graph stays connected, and the two
regions on the two sides of $e$, which are different because the cycle separates the plane (the Jordan curve theorem for polygons), merge into one: $E$ and $F$
both drop by $1$, and $V-E+F$ is unchanged.
\end{proof}

\begin{theorem}[Descartes]\label{geometry:thm:descartesdefect}
The angular defects of a convex polyhedron add up to $4\pi$.
\end{theorem}

\begin{proof}
A face with $n$ sides has angles adding up to $(n-2)\pi$. Summing over the faces, and using that each edge borders two faces, the sum of all face angles is
$\sum_f(n_f-2)\pi=(2E-2F)\pi$. Hence
\[
\sum_v\delta_v=2\pi V-(2E-2F)\pi=2\pi(V-E+F)=4\pi
\]
by Euler's formula.
\end{proof}

The theorem appears in a manuscript of Descartes from around 1630, \emph{De solidorum elementis}, known only through a copy made by Leibniz in 1676 and rediscovered
in the nineteenth century. It is Gauss--Bonnet for polyhedra: the curvature of a polyhedron is concentrated at its vertices, each carrying the amount $\delta_v$,
and the total is $4\pi$, as for every egg. Approximating a smooth convex surface by polyhedra, the defects become the curvature integrals of
Theorem~\ref{geometry:thm:gbconvex}. And so the book ends where geometry began, with the angles of a polygon, now seen as a shadow of the topology of the sphere.

\section{Exercises}

\begin{exercise}[$\star$]\label{geometry:ex:graphcurv}
Prove the formula $\kappa=\frac{x'y''-y'x''}{(x'^2+y'^2)^{3/2}}$ for any regular parametrisation, and deduce $\kappa=\frac{f''}{(1+f'^2)^{3/2}}$ for the graph $y=f(x)$.
\end{exercise}

\begin{exercise}[$\star$]\label{geometry:ex:ellipsevertex}
Show that the ellipse $x=a\cos t$, $y=b\sin t$ ($a>b>0$) has curvature $\kappa=\frac{ab}{(a^2\sin^2t+b^2\cos^2t)^{3/2}}$ and exactly four vertices.
\end{exercise}

\begin{exercise}[$\star$]\label{geometry:ex:constcurv}
Show that a curve with constant curvature $\kappa\ne0$ is an arc of a circle of radius $\frac1{|\kappa|}$.
\end{exercise}

\begin{exercise}[$\star$]\label{geometry:ex:totalcurv}
Compute the total curvature of the circle traversed once, and of the figure eight $x=\sin2t$, $y=\sin t$, $t\in[0,2\pi]$.
\end{exercise}

\begin{exercise}[$\star\star$]\label{geometry:ex:clothoid}
Describe the unit-speed curve with $\kappa(s)=s$ (the \emph{clothoid}, used to design road and railway bends). Show that it has a limit point as $s\to\infty$.
\end{exercise}

\begin{exercise}[$\star$]\label{geometry:ex:surfacesK}
Compute $K$ at the origin for $z=x^2+3y^2$, $z=x^2-y^2$, $z=x^2$, and for the monkey saddle $z=x^3-3xy^2$.
\end{exercise}

\begin{exercise}[$\star\star$]\label{geometry:ex:torusK}
Parametrise the torus by $\bigl((R+r\cos v)\cos u,(R+r\cos v)\sin u,r\sin v\bigr)$, $0<r<R$. Using that $K=\frac{\cos v}{r(R+r\cos v)}$ and
$dA=r(R+r\cos v)\,du\,dv$, compute $\iint K\,dA$. Where is $K$ positive, negative, zero?
\end{exercise}

\begin{exercise}[$\star$]\label{geometry:ex:platonic}
Check Descartes' theorem for the five Platonic solids.
\end{exercise}

\begin{exercise}[$\star\star$]\label{geometry:ex:fiveplatonic}
Using Euler's formula, show that a convex polyhedron whose faces are all $p$-gons, with $q$ faces at each vertex, has $\frac1p+\frac1q>\frac12$. Deduce that there
are at most five Platonic solids.
\end{exercise}

\begin{exercise}[$\star$]\label{geometry:ex:football}
A convex polyhedron has only pentagonal and hexagonal faces, with three faces at each vertex. Show that it has exactly $12$ pentagons.
\end{exercise}

\begin{exercise}[$\star\star$]\label{geometry:ex:nomap}
Let $A$ be the North Pole of the unit sphere, and $B$, $C$ two points on the circle of latitude at distance $a$ from $A$ (small), with $BC$ along a great circle of length
$a$. Let $M$ be the midpoint of the great circle arc $BC$. Compute $AM$ with the spherical law of cosines and compare with the value $\frac{\sqrt3}2a$ that a map
preserving the distances $AB$, $AC$, $BC$, $BM$, $CM$ would force in the plane. Conclude that no map from a piece of the sphere to the plane preserves all distances.
\end{exercise}

\section{Hints}

\begin{hint}{geometry:ex:graphcurv}
Use $ds=\sqrt{x'^2+y'^2}\,dt$ and $\tan\theta=\frac{y'}{x'}$.
\end{hint}
\begin{hint}{geometry:ex:ellipsevertex}
The curvature is extremal where $a^2\sin^2t+b^2\cos^2t$ is.
\end{hint}
\begin{hint}{geometry:ex:constcurv}
Use Theorem~\ref{geometry:thm:fundcurve}.
\end{hint}
\begin{hint}{geometry:ex:totalcurv}
Follow the direction of the tangent.
\end{hint}
\begin{hint}{geometry:ex:clothoid}
Write $\gamma(s)$ as an integral of $\bigl(\cos\frac{u^2}2,\sin\frac{u^2}2\bigr)$: these are the Fresnel integrals.
\end{hint}
\begin{hint}{geometry:ex:surfacesK}
Compute the Hessians.
\end{hint}
\begin{hint}{geometry:ex:torusK}
The integral in $v$ of $\cos v$ vanishes.
\end{hint}
\begin{hint}{geometry:ex:platonic}
The defect at a vertex where $q$ regular $p$-gons meet is $2\pi-q\pi\frac{p-2}p$.
\end{hint}
\begin{hint}{geometry:ex:fiveplatonic}
Count: $pF=2E=qV$.
\end{hint}
\begin{hint}{geometry:ex:football}
Use Descartes: each vertex carries the same defect, or count with Euler.
\end{hint}
\begin{hint}{geometry:ex:nomap}
$AM\perp BC$, so the right-angled spherical triangle $AMB$ gives $\cos a=\cos(AM)\cos\frac a2$.
\end{hint}

\section{Solutions}

\begin{solution}{geometry:ex:graphcurv}
\idea Curvature is the rate of turning of the tangent per unit length.
\writeup With $\theta=\arctan\frac{y'}{x'}$ (locally, up to a constant multiple of $\pi$), $\frac{d\theta}{dt}=\frac{x'y''-y'x''}{x'^2+y'^2}$, and $\frac{ds}{dt}=\sqrt{x'^2+y'^2}$.
Dividing, $\kappa=\frac{d\theta}{ds}=\frac{x'y''-y'x''}{(x'^2+y'^2)^{3/2}}$. For $x=t$, $y=f(t)$: $x'=1$, $x''=0$, so $\kappa=\frac{f''}{(1+f'^2)^{3/2}}$.
\end{solution}

\begin{solution}{geometry:ex:ellipsevertex}
\idea Plug into the formula, then look at a single trigonometric function.
\writeup With $x'=-a\sin t$, $y'=b\cos t$, $x''=-a\cos t$, $y''=-b\sin t$: $x'y''-y'x''=ab\sin^2t+ab\cos^2t=ab$, and the formula follows. The denominator involves
$g(t)=a^2\sin^2t+b^2\cos^2t=b^2+(a^2-b^2)\sin^2t$, which is extremal exactly when $\sin2t=0$: $t=0,\frac\pi2,\pi,\frac{3\pi}2$. So $\kappa$ has maxima $\frac a{b^2}$ at
$t=0,\pi$ and minima $\frac b{a^2}$ at $t=\pm\frac\pi2$, and no other critical points: exactly four vertices.
\end{solution}

\begin{solution}{geometry:ex:constcurv}
\idea The circle is a solution; uniqueness does the rest.
\writeup The circle of radius $\frac1{|\kappa|}$, traversed counterclockwise if $\kappa>0$ and clockwise if $\kappa<0$, has curvature $\kappa$. By Theorem~\ref{geometry:thm:fundcurve},
every curve with this curvature is obtained from it by a rotation and a translation, so it is (an arc of) a circle of the same radius.
\end{solution}

\begin{solution}{geometry:ex:totalcurv}
\idea Count how many times the tangent turns.
\writeup For the circle the tangent makes one full turn: $2\pi$. For the figure eight, use symmetry: $x(t+\pi)=x(t)$ and $y(t+\pi)=-y(t)$, so the second loop
($t\in[\pi,2\pi]$) is the mirror image of the first in the $x$-axis, traced with the same speed. A reflection changes the sign of $x'y''-y'x''$, hence of the signed curvature,
so the integral of $\kappa\,ds$ over the second loop is minus the integral over the first. The total curvature is $0$: the tangent turns one way on one loop and back on the other.
\end{solution}

\begin{solution}{geometry:ex:clothoid}
\idea The curvature grows linearly, so the curve winds into a tighter and tighter spiral.
\writeup By Theorem~\ref{geometry:thm:fundcurve}, $\theta(s)=\frac{s^2}2$ and
\[\gamma(s)=\Bigl(\int_0^s\cos\frac{u^2}2\,du,\ \int_0^s\sin\frac{u^2}2\,du\Bigr).\]
The substitution $v=\frac{u^2}2$ gives
$\int\cos\frac{u^2}2du=\int\frac{\cos v}{\sqrt{2v}}dv$, which converges as $s\to\infty$ by the alternating series argument (the integrals over $[k\pi,(k+1)\pi]$ alternate in sign and
decrease in size). So $\gamma(s)$ has a limit, $\bigl(\frac{\sqrt\pi}2,\frac{\sqrt\pi}2\bigr)$ by the Fresnel integrals: the curve spirals into this point. Engineers use the
clothoid between a straight road ($\kappa=0$) and a circular bend, so that the steering wheel turns at a constant rate.
\end{solution}

\begin{solution}{geometry:ex:surfacesK}
\idea $K$ is the determinant of the Hessian at a point with horizontal tangent plane.
\writeup $z=x^2+3y^2$: $H=\operatorname{diag}(2,6)$, $K=12$. $z=x^2-y^2$: $H=\operatorname{diag}(2,-2)$, $K=-4$. $z=x^2$: $K=0$ (a parabolic cylinder). Monkey saddle:
all second derivatives vanish at $0$, so $K=0$ there, although the surface is far from flat: it goes up and down three times around the origin (room for the
monkey's tail).
\end{solution}

\begin{solution}{geometry:ex:torusK}
\idea Positive curvature outside, negative inside, and they cancel.
\writeup We have $K\,dA=\frac{\cos v}{r(R+r\cos v)}\,r(R+r\cos v)\,du\,dv=\cos v\,du\,dv$, so
\[\iint K\,dA=\int_0^{2\pi}du\int_0^{2\pi}\cos v\,dv=0=2\pi\chi(\text{torus}).\]
Then $K>0$ on the outer half
($|v|<\frac\pi2$, where the torus looks like an egg), $K<0$ on the inner half (a saddle), and $K=0$ on the top and bottom circles $v=\pm\frac\pi2$.
\end{solution}

\begin{solution}{geometry:ex:platonic}
\idea Defect per vertex times number of vertices.
\writeup Tetrahedron ($p=3$, $q=3$, $V=4$): defect $2\pi-\pi=\pi$, total $4\pi$. Cube ($4,3$, $V=8$): $\frac\pi2$, total $4\pi$. Octahedron ($3,4$, $V=6$): $2\pi-\frac{4\pi}3=\frac{2\pi}3$, total
$4\pi$. Dodecahedron ($5,3$, $V=20$): $2\pi-3\cdot\frac{3\pi}5=\frac\pi5$, total $4\pi$. Icosahedron ($3,5$, $V=12$): $2\pi-\frac{5\pi}3=\frac\pi3$, total $4\pi$.
\end{solution}

\begin{solution}{geometry:ex:fiveplatonic}
\idea Euler's formula leaves room only for small $p$ and $q$.
\writeup Counting pairs (face, edge) and (vertex, edge): $pF=2E=qV$. So $V-E+F=E\bigl(\frac2q-1+\frac2p\bigr)=2>0$, hence $\frac1p+\frac1q>\frac12$. With $p,q\ge3$ the
solutions are $(3,3)$, $(3,4)$, $(4,3)$, $(3,5)$, $(5,3)$: tetrahedron, octahedron, cube, icosahedron, dodecahedron. So there are at most five, and all five exist.
\end{solution}

\begin{solution}{geometry:ex:football}
\idea Hexagons are flat enough not to count.
\writeup Let $f_5$, $f_6$ be the numbers of pentagons and hexagons. Then $2E=5f_5+6f_6$, $3V=2E$ and $F=f_5+f_6$. Euler's formula gives
$\frac{5f_5+6f_6}3-\frac{5f_5+6f_6}2+f_5+f_6=2$, that is $f_5\bigl(\frac53-\frac52+1\bigr)+f_6\bigl(2-3+1\bigr)=2$, so $\frac{f_5}6=2$ and $f_5=12$, whatever $f_6$ is. A football
has $12$ pentagons and $20$ hexagons; the fullerene molecule C$_{60}$ has the same shape.
\end{solution}

\begin{solution}{geometry:ex:nomap}
\idea On the sphere, the median of an equilateral triangle is longer than in the plane.
\writeup By symmetry $AM$ is perpendicular to $BC$, and the right-angled spherical Pythagorean theorem in $AMB$ gives $\cos a=\cos(AM)\cos\frac a2$, so
$\cos(AM)=\frac{\cos a}{\cos(a/2)}$. Expanding, $\frac{1-a^2/2+a^4/24}{1-a^2/8+a^4/384}=1-\frac38a^2-\frac1{128}a^4+O(a^6)$, while the planar value $\frac{\sqrt3}2a$ has cosine
$1-\frac38a^2+\frac{3}{128}a^4+O(a^6)$. The exact values differ for every $0<a<\frac\pi2$: for instance, for $a=0.5$, $AM\approx0.43767$ on the sphere and
$\frac{\sqrt3}2a\approx0.43301$ in the plane. A map preserving the five distances $AB=AC=BC=a$, $BM=CM=\frac a2$ would send $B,M,C$ to collinear points (since
$BM+MC=BC$ in the plane forces $M$ onto the segment) and $A$ to a point at distance $a$ from $B$ and $C$, so it would force $AM=\frac{\sqrt3}2a$. Since this is false on the
sphere, no map from a neighbourhood of $A$ to the plane preserves all distances.
\end{solution}

\section*{Sources and further reading}

\begin{itemize}
\item C.~F.~Gauss, \emph{Disquisitiones generales circa superficies curvas}, G\"ottingen, 1827 (the Theorema Egregium, paraphrased in the epigraph).
\item M.~P.~do Carmo, \emph{Differential Geometry of Curves and Surfaces}, Prentice-Hall, 1976 (Theorema Egregium, Gauss--Bonnet, the four vertex theorem).
\item S.~Mukhopadhyaya, \emph{New methods in the geometry of a plane arc}, Bull.\ Calcutta Math.\ Soc.\ 1 (1909); A.~Kneser, \emph{Bemerkungen \"uber die Anzahl der
Extreme der Kr\"ummung auf geschlossenen Kurven}, 1912; R.~Osserman, \emph{The four-or-more vertex theorem}, American Mathematical Monthly 92 (1985), 332--337.
\item H.~Hopf, \emph{\"Uber die Drehung der Tangenten und Sehnen ebener Kurven}, Compositio Mathematica 2 (1935).
\item P.~J.~Federico, \emph{Descartes on Polyhedra: A Study of the De Solidorum Elementis}, Springer, 1982 (Descartes' manuscript and Leibniz's copy).
\item D.~S.~Richeson, \emph{Euler's Gem: The Polyhedron Formula and the Birth of Topology}, Princeton University Press, 2008.
\end{itemize}
